\subsection{The coordinate completion on a general forest}
\label{sec:general-forest}

The paired argument used finite support and ancestral cuts, rather than
finite component heights. We record its resulting scope, while keeping
the distinction between a coordinate completion and a full dual explicit.
Let $T$ be any countably infinite rooted forest in which every vertex
has finitely many ancestors. Define $\rho_T$ by \eqref{eq:segment-norm}
with $T$ in place of $\mathcal F$, let $J(T)$ be the completion, and put
\[
 X(T)=\overline{\operatorname{span}}^{\|\cdot\|}\{e_v^*:v\in T\}
 \subset J(T)^*.
\]
Thus $X(T)$ is the completion of finite coordinate vectors in the polar
norm. We do not identify it with all of $J(T)^*$.

A finite ancestral set cuts each segment into a prefix and suffix.
The projections $P_H,Q_H$ are therefore contractive on $J(T)$ and
$X(T)$, and finite vectors outside $H$ are dense in $Q_HX(T)$.
Every finite set has finite ancestral closure, so an increasing finite
ancestral exhaustion gives $P_Hx\to x$ for all $x\in X(T)$.
The proof of Lemma~\ref{lem:paired} applies unchanged to these objects:
there are finitely many nonzero masses $m_g$, their initial segments
are disjoint, and unattached roots give no crossing term. These are
all the properties used in the test construction, the simultaneous
suffix cuts, and the passage to completed tails. Consequently
\eqref{eq:paired} and \eqref{eq:paired-two-parameters} hold in $X(T)$,
and the same arbitrary-center approximation gives
\begin{equation}\label{eq:general-forest-modulus}
 \widehat\delta_{X(T)}(t)\ge\sqrt{1+t^2/12}-1\qquad(t>0).
\end{equation}
This also explains why reflexivity is used for the example's structural
identification, rather than as a hypothesis of the paired estimate.

For $T=\mathbb N^{<\omega}$, the coordinate completion $X(T)$ is
the countably branching James-tree predual $B_\infty$, as described
in Banakh~\cite[\S2.II]{Banakh2000} following
Ghoussoub--Maurey--Schachermayer. We give its path obstruction in
the displayed norm. Each root-path
functional $p_v=\sum_{w\preceq v}e_w^*$ again has norm one. The
antichain calculation in \eqref{eq:antichain-isometry} uses only
disjoint segments and applies to $X(T)$ as well. Thus, for any
$\phi\in X(T)^*$ and finitely many children of a fixed vertex,
\[
 \left|\sum_j a_j\phi(e_{v^\frown j}^*)\right|
 \le\|\phi\|\left(\sum_j a_j^2\right)^{1/2}.
\]
Taking $a_j=\phi(e_{v^\frown j}^*)$ shows that these evaluations are
square summable, and hence tend to zero. The child increments are
therefore weakly null against the whole continuous dual. Applying
Lemma~\ref{lem:bounded-tree} with $m=M=c=1$ proves
$\overline\delta_N(\alpha/(2\beta))=0$ for every equivalent norm
$\alpha\|\cdot\|\le N\le\beta\|\cdot\|$ on this coordinate
completion. This argument does not require it to be reflexive.
